\section{Agent Integration}
\label{sec:agent-integration}

AI agents are first-class users of Timeline Compiler. They are not an afterthought wired to
a human-facing CLI; they are a primary consumer of the IR, the primary beneficiary of a
published JSON Schema, and the primary audience for the MCP server surface. This section
designs the full agent integration path: from raw inputs through ingestion to a valid IR,
through validation and repair, and through round-trip editing without loss of provenance.

The fundamental unit of agent work is the \emph{ingestion task}: given a \textbf{data source}
(work items, prose, structured export) and a \textbf{natural-language prompt}, produce a valid
Timeline IR. The prompt is a first-class input---it determines what is selected, how it is
framed, which entities are promoted to milestones, and which are dropped entirely. An agent
that ignores the prompt and imports everything it can find is not doing its job.


%%============================================================================
\subsection{The Ingestion Contract}
\label{sec:ingestion-contract}
%%============================================================================

The ingestion contract defines what the ingestion path is permitted to do with the inputs it
receives. Table~\ref{tab:ingestion-contract} captures the four categories of ingestion
decision. Violations of this contract---particularly silent data loss or silent inference---
are the most common source of incorrect or misleading timelines.

\begin{table}[h]
\centering
\small
\begin{tabular}{lp{4.5cm}p{6cm}}
\toprule
\textbf{Category} & \textbf{Definition} & \textbf{Examples} \\
\midrule
\textbf{Assumed} & Facts taken from the prompt or schema without source evidence & Title derived from prompt instruction; version \texttt{"1.0"} \\
\addlinespace
\textbf{Inferred} & Derived from source data by a deterministic rule & \texttt{axis\_unit} from date span; \texttt{time\_range} from min/max activity dates; track from \texttt{AreaPath} \\
\addlinespace
\textbf{Defaulted} & Omitted from IR, relying on documented schema default & \texttt{status: planned} omitted when all items are future; \texttt{locale} omitted \\
\addlinespace
\textbf{Rejected} & Source data discarded, reason logged & Bugs dropped when prompt says ``features and milestones only''; items outside prompted time window; duplicate or zero-duration events with no label \\
\bottomrule
\end{tabular}
\caption{Ingestion decision categories}
\label{tab:ingestion-contract}
\end{table}

\subsubsection{The Prompt as First-Class Input}
\label{sec:prompt-as-input}

The prompt directs every non-trivial ingestion decision. Table~\ref{tab:prompt-controls}
lists what the prompt controls and what the ingestion layer must extract from it before
touching the source data.

\begin{table}[h]
\centering
\small
\begin{tabular}{lp{9cm}}
\toprule
\textbf{Prompt Signal} & \textbf{Ingestion Effect} \\
\midrule
Time horizon (\textit{``Q1--Q3 2026''}) & Sets or constrains \texttt{time\_range}; items outside window are rejected \\
Entity filter (\textit{``features and milestones only''}) & Restricts work item types included; bugs, tasks dropped \\
Track grouping (\textit{``group by team''}) & Maps source field (e.g., \texttt{AreaPath}) to tracks \\
Emphasis (\textit{``highlight at-risk items''}) & Promotes \texttt{status} field; may add annotation \\
Framing (\textit{``executive summary''}) & Reduces label verbosity; promotes high-level items \\
Title / subtitle & Populates \texttt{metadata.title} / \texttt{metadata.subtitle} \\
Milestone promotion (\textit{``mark GA as a milestone''}) & Converts matching activity to milestone entity \\
\bottomrule
\end{tabular}
\caption{Prompt signals and their ingestion effects}
\label{tab:prompt-controls}
\end{table}

\subsubsection{What Ingestion May Not Do}

\begin{itemize}
  \item \textbf{Must not compute dates.} If a source item has no date, it must use
        \texttt{tbd}, not invent one. Timeline Grammar does not compute schedules
        (see \S\ref{sec:scope}).
  \item \textbf{Must not collapse ambiguous status silently.} An ADO work item in state
        ``On Hold'' that has no clear IR equivalent must log a rejection warning, not
        silently map to \texttt{planned}.
  \item \textbf{Must not import the whole backlog.} The prompt defines the editorial
        scope. Importing everything that exists and hoping the renderer trims it is a
        contract violation.
  \item \textbf{Must not generate non-stable IDs.} IDs derived from sequential numbers
        (\texttt{act-1}, \texttt{act-2}) are permitted only as a last resort. IDs should
        be derived from the item's title slug for git-friendliness (see
        \S\ref{sec:round-trip}).
\end{itemize}


%%============================================================================
\subsection{Ingestion Workflows}
\label{sec:ingestion-workflows}
%%============================================================================

The following subsections describe four concrete ingestion flows, each with a source
excerpt, a prompt, the resulting IR YAML, and an explicit account of what the prompt
selected, inferred, and rejected.

%% ---------------------------------------------------------------------------
\subsubsection{Azure DevOps Work Items}
\label{sec:ingest-ado}
%% ---------------------------------------------------------------------------

\textbf{Source format.} The ADO Work Items REST API \cite{ado-workitems} returns JSON
objects. The fields relevant to timeline ingestion are listed in
Table~\ref{tab:ado-field-mapping}.

\begin{table}[h]
\centering
\small
\begin{tabular}{lll}
\toprule
\textbf{ADO Field} & \textbf{IR Target} & \textbf{Notes} \\
\midrule
\texttt{System.Title} & \texttt{activity.label} / \texttt{milestone.label} & Verbatim or
  prompt-shortened \\
\texttt{System.State} & \texttt{activity.status} & Via Table~\ref{tab:ado-state-map} \\
\texttt{System.WorkItemType} & \texttt{activity.category} / entity type & Feature $\to$
  activity; Milestone $\to$ milestone \\
\texttt{System.IterationPath} & \texttt{span} or \texttt{start}/\texttt{end} & E.g.
  \texttt{Platform\textbackslash 2026\textbackslash Q2} $\to$ \texttt{span: 2026-Q2} \\
\texttt{System.AreaPath} & \texttt{track} & Leaf segment becomes track label/id \\
\texttt{Microsoft.VSTS.Scheduling.TargetDate} & \texttt{milestone.date} & Used when
  promoting to milestone \\
\texttt{System.Tags} & \texttt{activity.tags} & Passed through \\
\texttt{System.Id} & \texttt{metadata.ado\_id} & Provenance (not rendered) \\
\bottomrule
\end{tabular}
\caption{ADO work item field mapping}
\label{tab:ado-field-mapping}
\end{table}

\begin{table}[h]
\centering
\small
\begin{tabular}{ll}
\toprule
\textbf{ADO State} & \textbf{IR Status} \\
\midrule
New          & \texttt{planned} \\
Active       & \texttt{in-progress} \\
Resolved     & \texttt{done} \\
Closed       & \texttt{done} \\
Removed      & \texttt{cancelled} \\
On Hold      & \texttt{blocked} \\
\textit{(other)} & \texttt{planned} (warn: unmapped state) \\
\bottomrule
\end{tabular}
\caption{ADO state to IR status mapping}
\label{tab:ado-state-map}
\end{table}

\textbf{Source excerpt.}

\begin{lstlisting}[language=yaml,caption={ADO work item JSON (abridged)}]
# ADO API response (workItems array, fields selected)
- id: 1024
  fields:
    System.Title: "API Gateway Migration"
    System.State: "Active"
    System.WorkItemType: "Feature"
    System.IterationPath: "Platform\\2026\\Q2"
    System.AreaPath: "Platform Team"
- id: 1025
  fields:
    System.Title: "Mobile App 3.0 GA"
    System.State: "New"
    System.WorkItemType: "Milestone"
    System.AreaPath: "Mobile Team"
    Microsoft.VSTS.Scheduling.TargetDate: "2026-09-15"
- id: 1031
  fields:
    System.Title: "Fix login timeout bug"
    System.State: "Active"
    System.WorkItemType: "Bug"
    System.AreaPath: "Platform Team"
\end{lstlisting}

\textbf{Prompt.} \textit{``Show the Q1--Q3 2026 Platform and Mobile roadmap. Include
features and milestones only. Group by team. Drop bugs and tasks.''}

\textbf{Resulting IR.}

\begin{lstlisting}[language=yaml,caption={IR produced from ADO work items}]
version: "1.0"
metadata:
  title: Platform & Mobile Roadmap
  subtitle: Q1-Q3 2026
  time_range:
    start: 2026-Q1
    end: 2026-Q3
  axis_unit: quarter   # inferred: 3-quarter window -> quarter granularity

tracks:
  - id: platform
    label: Platform Team
  - id: mobile
    label: Mobile Team

activities:
  - id: api-gateway-migration
    label: API Gateway Migration
    track: platform
    span: 2026-Q2        # inferred from IterationPath "Platform\2026\Q2"
    status: in-progress   # mapped from ADO state "Active"
    category: feature
    metadata:
      source: ado
      ado_id: 1024
      ado_iteration: "Platform\\2026\\Q2"

milestones:
  - id: mobile-app-3-ga
    label: Mobile App 3.0 GA
    date: 2026-09-15    # from TargetDate field
    track: mobile
    status: planned      # mapped from ADO state "New"
    metadata:
      source: ado
      ado_id: 1025
\end{lstlisting}

\textbf{What the prompt did.}
\begin{itemize}
  \item \textbf{Selected:} Items 1024 (Feature) and 1025 (Milestone). Item 1031 (Bug)
        was rejected because the prompt said ``features and milestones only''.
  \item \textbf{Framed:} \texttt{time\_range} set to Q1--Q3 2026 from prompt instruction.
        Item 1025's target date (2026-09-15) falls within Q3, so it is included.
  \item \textbf{Inferred:} Tracks derived from \texttt{AreaPath} leaf segments. \texttt{axis\_unit}
        inferred as \texttt{quarter} from the three-quarter window.
  \item \textbf{Dropped:} Bug work item 1031 (not in filter). Had the time range excluded
        any item, it would be dropped with a logged rejection.
\end{itemize}


%% ---------------------------------------------------------------------------
\subsubsection{Natural Language Prose}
\label{sec:ingest-prose}
%% ---------------------------------------------------------------------------

When the source is free-form text---a planning document, a briefing note, an email
thread---the agent must extract temporal structure from prose. This is the most
ambiguity-prone ingestion path; the contract requires explicit handling of every point
of uncertainty.

\textbf{Source text.}

\begin{lstlisting}[caption={Natural language source}]
We are planning to migrate the data warehouse in the first half of 2026.
The new customer portal is targeted for Q3. Both initiatives are owned
by the Engineering org. An executive review is scheduled for 15 July 2026.
The legacy decommission date is TBD, pending the warehouse cutover.
\end{lstlisting}

\textbf{Prompt.} \textit{``Show this as a two-track roadmap: Infrastructure and Product.
Mark the executive review as a milestone. Include the decommission even though its date
is not fixed.''}

\textbf{Resulting IR.}

\begin{lstlisting}[language=yaml,caption={IR produced from natural language prose}]
version: "1.0"
metadata:
  title: 2026 Engineering Roadmap
  time_range:
    start: 2026-H1
    end: 2026-Q3
  axis_unit: quarter  # inferred: H1+Q3 window, quarterly detail

tracks:
  - id: infrastructure
    label: Infrastructure
  - id: product
    label: Product

activities:
  - id: data-warehouse-migration
    label: Data Warehouse Migration
    track: infrastructure
    span: 2026-H1
    status: planned

  - id: customer-portal
    label: New Customer Portal
    track: product
    span: 2026-Q3
    status: planned

  - id: legacy-decommission
    label: Legacy Decommission
    track: infrastructure
    start: tbd          # explicit: date not yet determined
    status: planned

milestones:
  - id: executive-review
    label: Executive Review
    date: 2026-07-15
    status: planned
\end{lstlisting}

\textbf{What the prompt did.}
\begin{itemize}
  \item \textbf{Structured:} Track assignments (Infrastructure vs.\ Product) came from the
        prompt. The source text only named ``Engineering org''.
  \item \textbf{Promoted:} The executive review was elevated to a milestone (not an
        activity) by explicit prompt instruction.
  \item \textbf{Included with uncertainty:} Legacy decommission included despite no
        fixed date; \texttt{start: tbd} preserves the fact on the timeline without
        inventing a date.
  \item \textbf{Inferred:} \texttt{time\_range} derived from the earliest (\texttt{2026-H1}
        start) and latest (\texttt{2026-07-15}) concrete dates. \texttt{axis\_unit}
        set to \texttt{quarter} as the coarsest granularity that fits the items.
\end{itemize}

\textbf{Ambiguity resolution.} ``First half of 2026'' is unambiguously \texttt{2026-H1}
(ISO half, per the date model). ``Q3'' without a year is inferred as \texttt{2026-Q3}
from context (all other dates are 2026). If the year were ambiguous, the agent must
surface a warning rather than guess.


%% ---------------------------------------------------------------------------
\subsubsection{GitHub Projects}
\label{sec:ingest-github}
%% ---------------------------------------------------------------------------

GitHub Projects (ProjectV2) \cite{github-projects,github-graphql-projectv2} supports
custom field types that carry roadmap-relevant data. The ingestion mapping is
summarised in Table~\ref{tab:github-field-mapping}.

\begin{table}[h]
\centering
\small
\begin{tabular}{lll}
\toprule
\textbf{GitHub ProjectV2 Field} & \textbf{IR Target} & \textbf{Notes} \\
\midrule
\texttt{title} (built-in) & \texttt{activity.label} / \texttt{milestone.label} & Verbatim or shortened \\
\texttt{status} (select field) & \texttt{activity.status} & Label mapped to IR enum \\
\texttt{DATE} custom field & \texttt{start} / \texttt{milestone.date} & Named field determines role \\
\texttt{ITERATION} custom field & \texttt{span} & Iteration title maps to quarter \\
\texttt{labels} & \texttt{activity.category} (first) / \texttt{activity.tags} & First label to category \\
\texttt{milestone} (linked) & \texttt{milestones} entry & Promoted to IR milestone \\
\texttt{repository} & \texttt{track} & When prompt says ``group by repo'' \\
\texttt{assignees} & \texttt{metadata.assignees} & Not rendered; preserved for provenance \\
\texttt{number} (issue \#) & \texttt{metadata.github\_issue} & Provenance link \\
\bottomrule
\end{tabular}
\caption{GitHub ProjectV2 field mapping}
\label{tab:github-field-mapping}
\end{table}

The ITERATION field type is the most useful for quarterly roadmaps: iteration titles
like ``2026 Q2'' map directly to \texttt{span: 2026-Q2}, making the ingestion path
from GitHub Projects to a quarterly roadmap nearly lossless.

\textbf{Gap note.} GitHub Projects does not have a first-class concept analogous to
``work item type'' (Feature vs.\ Bug vs.\ Task). Filtering by type requires label or
custom field conventions, which vary per team. The prompt must specify the filter
criterion explicitly.


%% ---------------------------------------------------------------------------
\subsubsection{Mermaid Timeline Syntax}
\label{sec:ingest-mermaid}
%% ---------------------------------------------------------------------------

Mermaid's \texttt{timeline} diagram \cite{mermaiddocs2024,mermaid2023} is the most
common structured timeline format already embedded in team documentation. Its syntax
maps cleanly to IR sections and activities.

\textbf{Source.}
\begin{lstlisting}[caption={Mermaid timeline source}]
timeline
    title Platform Evolution 2026
    section H1 2026
        API Gateway : Q1 : Migration complete
        Auth Refresh : Q2 : New OAuth2 provider
    section H2 2026
        Mobile 3.0 GA : Q3
        Legacy Cutover : Q4 : Final decommission
\end{lstlisting}

\textbf{Resulting IR.}
\begin{lstlisting}[language=yaml,caption={IR produced from Mermaid timeline syntax}]
version: "1.0"
metadata:
  title: Platform Evolution 2026
  time_range:
    start: 2026-Q1
    end: 2026-Q4
  axis_unit: quarter

tracks:
  - id: platform
    label: Platform      # default single-track when Mermaid has no swimlanes

sections:
  - id: h1-2026
    label: H1 2026
    time_range:
      start: 2026-Q1
      end: 2026-Q2

  - id: h2-2026
    label: H2 2026
    time_range:
      start: 2026-Q3
      end: 2026-Q4

activities:
  - id: api-gateway
    label: API Gateway
    track: platform
    span: 2026-Q1
    description: "Migration complete"
    status: planned

  - id: auth-refresh
    label: Auth Refresh
    track: platform
    span: 2026-Q2
    description: "New OAuth2 provider"
    status: planned

  - id: mobile-3-ga
    label: Mobile 3.0 GA
    track: platform
    span: 2026-Q3
    status: planned

  - id: legacy-cutover
    label: Legacy Cutover
    track: platform
    span: 2026-Q4
    description: "Final decommission"
    status: planned
\end{lstlisting}

\textbf{Mapping notes.}
Mermaid \texttt{timeline} does not have a swimlane concept; all events are rendered in
a single horizontal flow grouped by section. When promoting to IR, a default single track
is created unless the prompt specifies a multi-track structure. Mermaid's colon-separated
description field maps to \texttt{description}. Because Mermaid \texttt{gantt}
\cite{mermaidgantt2024} uses a different syntax (task dependencies, progress), Mermaid
gantt ingestion requires a separate mapping pass; that mapping is out of scope for the
initial IR but is an extension point (see \S\ref{sec:scope}).


%%============================================================================
\subsection{Reliable IR Generation by Agents}
\label{sec:agent-ir-generation}
%%============================================================================

LLM-based agents generate IR reliably when three conditions hold: (1) the target schema
is machine-readable and published, (2) the generation call is constrained to valid
output shapes, and (3) the IR's structural choices minimise the surface area for
generation errors.

%% ---------------------------------------------------------------------------
\subsubsection{JSON Schema Publication}
\label{sec:json-schema}
%% ---------------------------------------------------------------------------

The IR schema must be published as a JSON Schema document \cite{openai-structured-outputs}.
This is a binding requirement established by the positioning analysis
(\S\ref{sec:comparison}): without a machine-readable schema, agent generation is
unreliable. The schema is a first-class deliverable of this project, versioned alongside
the spec.

The top-level JSON Schema structure mirrors the IR document structure directly:

\begin{lstlisting}[language=yaml,caption={Abbreviated JSON Schema for the IR (YAML representation)}]
# $schema: https://json-schema.org/draft/2020-12/schema
# $id: https://timeline-compiler.dev/schema/v1/timeline.json
title: Timeline IR
type: object
required: [version, metadata, tracks, activities]
properties:
  version:
    type: string
    description: "Specification version, e.g. \"1.0\""

  metadata:
    type: object
    required: [title, time_range]
    properties:
      title:       { type: string }
      subtitle:    { type: string }
      time_range:
        type: object
        required: [start]
        properties:
          start: { "$ref": "#/$defs/Date" }
          end:   { "$ref": "#/$defs/Date" }
      axis_unit:   { "$ref": "#/$defs/AxisUnit" }
      theme:       { type: string }
      locale:      { type: string }

  tracks:
    type: array
    minItems: 1
    items: { "$ref": "#/$defs/Track" }

  activities:
    type: array
    items: { "$ref": "#/$defs/Activity" }

  milestones:
    type: array
    items: { "$ref": "#/$defs/Milestone" }

  # groups, annotations, sections, legend all optional arrays

$defs:
  ID:
    type: string
    pattern: "^[a-z][a-z0-9-]*$"

  Date:
    type: string
    description: "ISO date, quarter (2026-Q2), half (2026-H1), relative (+3m),
                  symbolic (now), or uncertain token (tbd, ongoing, unknown,
                  or tilde-prefix ~2026-Q2)"

  AxisUnit:
    type: string
    enum: [day, week, month, quarter, half, year]

  Status:
    type: string
    enum: [planned, in-progress, done, at-risk, blocked, cancelled, tentative]

  Track:
    type: object
    required: [id, label]
    properties:
      id:    { "$ref": "#/$defs/ID" }
      label: { type: string }
      color: { type: string }
      order: { type: integer, minimum: 0 }

  Activity:
    type: object
    required: [id, label, track]
    oneOf:
      - required: [start]    # start + optional end
      - required: [span]     # span shorthand
    properties:
      id:       { "$ref": "#/$defs/ID" }
      label:    { type: string }
      track:    { type: string }   # bare string ref to Track.id
      start:    { "$ref": "#/$defs/Date" }
      end:      { "$ref": "#/$defs/Date" }
      span:     { "$ref": "#/$defs/Date" }
      status:   { "$ref": "#/$defs/Status", default: planned }
      progress: { type: number, minimum: 0, maximum: 1 }
      category: { type: string }
      group:    { type: string }
      description: { type: string }
      metadata: { type: object }

  Milestone:
    type: object
    required: [id, label, date]
    properties:
      id:     { "$ref": "#/$defs/ID" }
      label:  { type: string }
      date:   { "$ref": "#/$defs/Date" }
      track:  { type: string }
      status: { "$ref": "#/$defs/Status", default: planned }
      category: { type: string }
      icon:     { type: string }
      description: { type: string }
      metadata: { type: object }
\end{lstlisting}

The schema is published at a versioned URL. Minor version bumps to the IR may add
optional fields; the schema handles this via \texttt{additionalProperties: true} for
\texttt{metadata} maps and by allowing unknown top-level fields on minor revisions.

\paragraph{New cite key needed.} The JSON Schema specification itself does not have a
cite key in \texttt{references.bib}. A key \texttt{json-schema2020} pointing to the
IETF JSON Schema draft (2020-12) should be added by David.


%% ---------------------------------------------------------------------------
\subsubsection{Structured Output Constraints}
\label{sec:structured-outputs}
%% ---------------------------------------------------------------------------

Platforms that support constrained generation \cite{openai-structured-outputs} should
use the published JSON Schema as the output schema for the IR generation call. This
eliminates entire classes of generation errors:

\begin{itemize}
  \item Required fields can never be missing when the schema enforces \texttt{required}.
  \item ID format is enforced by the \texttt{pattern} constraint on the \texttt{ID}
        definition, preventing spaces, uppercase, or other invalid characters.
  \item Status values are constrained to the enum; the model cannot hallucinate
        ``completed'' or ``active''.
  \item References are just bare strings---the model emits \texttt{track: "platform"},
        not a nested object, which is trivially correct.
\end{itemize}

When structured output is not available, the agent prompt must include the JSON Schema
inline and instruct the model to validate its own output before returning it.

%% ---------------------------------------------------------------------------
\subsubsection{IR Design Choices That Help Agents}
\label{sec:ir-agent-friendly}
%% ---------------------------------------------------------------------------

Several deliberate IR design choices reduce the error rate for LLM-generated documents:

\begin{description}
  \item[Optional fields default to safe values.] \texttt{status} defaults to
        \texttt{planned}; \texttt{progress} defaults to \texttt{null}. An agent that
        omits these fields produces a valid document. Agents only need to emit fields
        they have evidence for.

  \item[\texttt{span} shorthand reduces reasoning load.] Instead of requiring an agent
        to compute both a \texttt{start} and \texttt{end} date for a quarter-length
        activity, it can emit \texttt{span: 2026-Q2} and let the compiler expand it.
        This avoids off-by-one boundary errors (is Q2 end June 30 or July 1?).

  \item[References are bare strings.] \texttt{track: platform} is far easier for an
        LLM to emit correctly than a nested reference object. The schema context (field
        name) determines the target type; no wrapper object is needed.

  \item[IDs are kebab-case slugs.] Agents derive IDs from labels:
        ``API Gateway Migration'' $\to$ \texttt{api-gateway-migration}. This is a
        simple and reliable transformation. UUIDs would be safe but opaque; sequential
        numbers would collide with re-generation.

  \item[No dependency scheduling.] Since the IR is a Timeline Grammar (not a Task
        Scheduling Grammar), there is no dependency resolution. Each activity is
        independently valid. Agents do not need to compute a topological order or
        check resource constraints.
\end{description}


%%============================================================================
\subsection{Validation Pipeline}
\label{sec:validation-pipeline}
%%============================================================================

Validation is layered. Each layer is a gate: failure at an earlier layer produces
errors that block later layers. This design prevents cryptic failures (e.g., a
reference-resolution error caused by a silently malformed ID).

\begin{table}[h]
\centering
\small
\begin{tabular}{clp{4cm}p{4.5cm}}
\toprule
\textbf{Layer} & \textbf{Name} & \textbf{What it checks} & \textbf{Failure mode} \\
\midrule
1 & Syntactic parse       & Valid YAML or JSON; no parse errors & Hard error: stop \\
2 & Schema conformance    & Required fields present; types match; enum values valid; ID regex; \texttt{oneOf} for start/span & Hard error: list all violations \\
3 & Well-formedness       & Referential integrity; ID uniqueness; date ordering; progress bounds; group acyclicity & Hard error: list all violations \\
4 & Render-readiness      & See \S\ref{sec:rendering} (Barbara) & Error or warning depending on severity \\
5 & Semantic advisory     & Degenerate documents; empty \texttt{activities} and \texttt{milestones}; items outside \texttt{time\_range} & Warning only \\
\bottomrule
\end{tabular}
\caption{Validation pipeline layers}
\label{tab:validation-layers}
\end{table}

\textbf{Layer 4 dependency.} Render-readiness validation---checking that the IR can
produce a legible visual output---is defined and owned by Barbara in \S\ref{sec:rendering}.
This includes checks such as overlapping activities on the same track that cannot be
visually separated, labels that exceed track width at the chosen axis unit, and milestone
dates that fall outside the document's \texttt{time\_range}. The validation pipeline
calls Barbara's render-readiness checks as Layer 4; this section does not duplicate them.

\subsubsection{Errors vs.\ Warnings}

\begin{description}
  \item[Error] Prevents rendering. The document is invalid and the renderer must refuse
        to proceed. Errors come from Layers 1--3 and from Layer 4 severity ``fatal''.
        The validator must return \emph{all} errors in a single pass, not just the first.
  \item[Warning] The document is valid but the output may be suboptimal. Warnings come
        from Layer 4 (advisory) and Layer 5. Renderers proceed and may annotate the
        output (e.g., a clipped activity) or log the warning and continue.
\end{description}


%%============================================================================
\subsection{Error Handling and Agent Repair Cycle}
\label{sec:error-repair}
%%============================================================================

\subsubsection{Error Message Format}

Every validation error is \textbf{path-anchored}: it identifies the exact location in
the document, the nature of the violation, and a suggested fix. This design enables
an agent to mechanically apply repairs without re-reading the whole document.

The error message format is:

\begin{lstlisting}[caption={Validation error message format}]
ERROR  <path>: <description>
       Expected: <constraint>
       Found:    <actual value>
       Fix:      <suggested correction>
\end{lstlisting}

\textbf{Concrete error examples:}

\begin{lstlisting}[caption={Example validation errors with suggested fixes}]
# Missing required field
ERROR  activities[0].id: required field 'id' is missing
       Expected: string matching ^[a-z][a-z0-9-]*$
       Found:    (absent)
       Fix:      add id: api-gateway-migration

# Invalid ID format
ERROR  activities[1].id: "API Migration" does not match ^[a-z][a-z0-9-]*$
       Expected: kebab-case slug
       Found:    "API Migration"
       Fix:      id: api-migration

# Unresolved track reference
ERROR  activities[2].track: "eng" references a non-existent track
       Expected: one of [platform, mobile, data]
       Found:    "eng"
       Fix:      track: platform  (or add a track with id: eng)

# Invalid status value
ERROR  activities[3].status: "active" is not a valid Status
       Expected: one of [planned, in-progress, done, at-risk,
                          blocked, cancelled, tentative]
       Found:    "active"
       Fix:      status: in-progress

# Date ordering violation
ERROR  activities[4]: end precedes start
       Expected: end >= start
       Found:    start: 2026-Q3, end: 2026-Q1
       Fix:      swap dates, or use span: 2026-Q3 if this is
                 a single-quarter activity

# Progress out of range
ERROR  activities[5].progress: 1.4 is outside [0, 1]
       Expected: number in [0, 1]
       Found:    1.4
       Fix:      progress: 1.0

# Duplicate ID
ERROR  milestones[1].id: "platform-launch" is already used by activities[2].id
       Expected: globally unique identifier
       Found:    duplicate
       Fix:      id: platform-launch-ms  (or another unique slug)

# Missing date source (neither start nor span present)
ERROR  activities[6]: no temporal anchor defined
       Expected: exactly one of: (start [+ optional end]) or span
       Found:    neither start nor span present
       Fix:      add  span: 2026-Q2  or  start: 2026-04-01
\end{lstlisting}

\subsubsection{Agent Repair Cycle}
\label{sec:repair-cycle}

The validation-error-repair loop is the core of reliable agent IR generation. The
cycle is:

\begin{enumerate}
  \item \textbf{Generate.} Agent generates an IR document, optionally using structured
        output constraints (\S\ref{sec:structured-outputs}).
  \item \textbf{Validate.} The IR is submitted to the \texttt{validate\_timeline} MCP
        tool (\S\ref{sec:mcp-tools}). All layers are run; all errors collected.
  \item \textbf{Report.} Errors are returned to the agent as a structured list with
        path, description, and suggested fix.
  \item \textbf{Repair.} The agent applies targeted fixes. Because errors are
        path-anchored and include suggested fixes, repairs are surgical: the agent
        patches \texttt{activities[2].track} rather than regenerating the entire document.
  \item \textbf{Re-validate.} The patched document is re-submitted. This loop repeats
        until validation passes with no errors (warnings may remain).
  \item \textbf{Render.} Once validation passes, the agent calls \texttt{render\_timeline}
        to produce output.
\end{enumerate}

In practice, structured output constraints eliminate most Layer~2 errors on the first
generation pass. Layers~3 (referential integrity, date ordering) are the most common
source of residual errors for agents generating multi-track documents.


%%============================================================================
\subsection{MCP Server}
\label{sec:mcp-server}
%%============================================================================

The MCP server \cite{mcp-spec} exposes Timeline Compiler as a tool-callable service for
AI agents. It is a first-class distribution target, not an afterthought (see the
distribution architecture in \S\ref{sec:distribution}). The server hosts four tools.

%% ---------------------------------------------------------------------------
\subsubsection{Tool Contracts}
\label{sec:mcp-tools}
%% ---------------------------------------------------------------------------

\paragraph{\texttt{validate\_timeline}}

Runs the full validation pipeline (Layers 1--5) against a provided IR document.
Returns all errors and warnings in structured form.

\begin{lstlisting}[language=yaml,caption={\texttt{validate\_timeline} tool contract}]
tool: validate_timeline
description: >
  Validate a Timeline IR document. Runs syntactic, schema, well-formedness,
  render-readiness, and semantic advisory checks. Returns structured errors
  and warnings with path anchors and suggested fixes.

input:
  ir:
    type: object | string   # IR as parsed object or raw YAML/JSON string
    required: true
  stop_at_layer:
    type: integer           # optional; 1-5; default: 5 (run all layers)
    required: false

output:
  valid:     boolean        # true iff zero errors (warnings do not block)
  errors:
    - path:    string       # e.g. "activities[2].track"
      code:    string       # machine-readable error code, e.g. UNRESOLVED_REF
      message: string       # human-readable description
      fix:     string       # suggested correction
  warnings:
    - path:    string
      code:    string
      message: string
  layers_run: integer       # how many layers were executed before stopping
\end{lstlisting}

\paragraph{\texttt{render\_timeline}}

Renders a valid Timeline IR to a visual output format. Returns the rendered output
as a base64-encoded string or a file path (for local MCP deployments).

\begin{lstlisting}[language=yaml,caption={\texttt{render\_timeline} tool contract}]
tool: render_timeline
description: >
  Render a Timeline IR document to SVG, PNG, PDF, or PPTX. The IR must
  be valid (validate_timeline returns valid: true) before rendering.
  Rendering is deterministic: identical IR + theme + format always
  produces bit-identical output.

input:
  ir:
    type: object | string   # IR as parsed object or raw YAML/JSON string
    required: true
  format:
    type: string
    enum: [svg, png, pdf, pptx]
    default: svg
    required: false
  theme:
    type: string            # theme identifier, e.g. "default", "corporate"
    default: "default"
    required: false
  width:
    type: integer           # output width in pixels (PNG/PDF only)
    required: false

output:
  format:  string           # echo of requested format
  data:    string           # base64-encoded output bytes
  mime_type: string         # e.g. "image/svg+xml"
  warnings:                 # render-layer warnings (does not block output)
    - message: string
\end{lstlisting}

\paragraph{\texttt{describe\_schema}}

Returns the full JSON Schema for the IR and field-level documentation. Agents call
this to bootstrap IR generation or to recover from schema errors.

\begin{lstlisting}[language=yaml,caption={\texttt{describe\_schema} tool contract}]
tool: describe_schema
description: >
  Return the JSON Schema for the Timeline IR, plus human-readable
  documentation for each field. Use this to understand what fields are
  required, their types, allowed values, and defaults.

input:
  version:
    type: string            # IR version to describe; default: latest
    required: false
  entity:
    type: string            # optional: limit to one entity, e.g. "Activity"
    required: false

output:
  version:     string       # IR version described
  schema:      object       # full JSON Schema document
  docs:
    - entity:  string       # entity name
      fields:
        - name:        string
          type:        string
          required:    boolean
          default:     string | null
          description: string
          example:     string
\end{lstlisting}

\paragraph{\texttt{suggest\_time\_range}}

A convenience tool for agents that need to derive a sensible \texttt{time\_range} and
\texttt{axis\_unit} from a collection of candidate activities before composing the full
document. This is especially useful when ingesting data from sources that lack explicit
time windows.

\begin{lstlisting}[language=yaml,caption={\texttt{suggest\_time\_range} tool contract}]
tool: suggest_time_range
description: >
  Given a list of activity date hints (start, end, span values), suggest
  appropriate time_range and axis_unit values for the IR metadata. Does not
  require a complete IR document.

input:
  dates:
    type: array             # list of date strings, e.g. ["2026-Q1", "2026-09-15"]
    items: { type: string }
    required: true
  prefer_unit:
    type: string            # optional hint: "quarter", "month", etc.
    required: false

output:
  time_range:
    start: string
    end:   string
  axis_unit: string         # recommended AxisUnit
  rationale: string         # explanation of the choice
\end{lstlisting}

\subsubsection{MCP Server Deployment}

The MCP server is deployed in two modes:

\begin{description}
  \item[Local server] Runs as a subprocess alongside the agent environment. Started
        via the CLI (\texttt{timeline mcp-server --port 3000}) or as an npm/pip
        package. Suitable for development, CI, and agent runtimes on the same host.
  \item[Hosted endpoint] A cloud-deployed instance of the MCP server. Suitable for
        cloud-native agent runtimes where a local subprocess is impractical.
\end{description}

Both modes expose identical tool contracts. The local server operates on the local
filesystem (output file paths are local); the hosted server returns base64-encoded
output in all cases.


%%============================================================================
\subsection{Round-Trip and Iterative Editing}
\label{sec:round-trip}
%%============================================================================

A Timeline IR document lives in version control alongside the source data it was derived
from. Humans and agents edit it directly. Re-syncing against an updated source must not
silently overwrite editorial decisions made in the IR. This section designs the mechanisms
that make round-trip editing safe.

%% ---------------------------------------------------------------------------
\subsubsection{Source Provenance via Metadata}
\label{sec:provenance}
%% ---------------------------------------------------------------------------

Every entity produced by an ingestion workflow includes a \texttt{metadata} block that
records the source of record:

\begin{lstlisting}[language=yaml,caption={Provenance metadata block (ADO example)}]
activities:
  - id: api-gateway-migration
    label: API Gateway Migration
    track: platform
    span: 2026-Q2
    status: in-progress
    metadata:
      source: ado
      ado_id: 1024
      ado_revision: 42           # ADO ETag for change detection
      ado_area: "Platform Team"
      ado_iteration: "Platform\\2026\\Q2"
      ingested_at: 2026-06-09
\end{lstlisting}

\begin{lstlisting}[language=yaml,caption={Provenance metadata block (GitHub Projects example)}]
activities:
  - id: customer-portal-redesign
    label: Customer Portal Redesign
    track: product
    span: 2026-Q3
    status: planned
    metadata:
      source: github
      github_project_id: PVT_kwDOA12345
      github_issue: 214
      github_url: "https://github.com/acme/portal/issues/214"
      ingested_at: 2026-06-09
\end{lstlisting}

The \texttt{source} key in \texttt{metadata} is a reserved provenance marker. The
ingester writes it; the renderer ignores it. Re-sync logic uses it to correlate IR
entities back to their source records.

\subsubsection{Re-Sync Strategy}
\label{sec:resync}

When source data is updated (e.g., a work item's state changes from Active to Resolved),
the re-sync operation must:

\begin{enumerate}
  \item \textbf{Match by source ID.} Use \texttt{metadata.ado\_id} or
        \texttt{metadata.github\_issue} as the stable foreign key, not the IR \texttt{id}.
        The IR \texttt{id} is a human-edited slug that may have been renamed.
  \item \textbf{Update only source-mapped fields.} Fields that the ingestion workflow
        maps from the source (e.g., \texttt{status}, \texttt{span} derived from
        IterationPath) are overwritten. Fields the human may have edited in the IR
        (\texttt{label}, \texttt{category}, \texttt{description}, \texttt{color},
        \texttt{progress}) are \emph{not} overwritten.
  \item \textbf{Preserve the IR \texttt{id}.} The kebab-case slug must not be
        regenerated on re-sync. Once set, it is stable.
  \item \textbf{Log what changed.} The re-sync operation must produce a structured
        change log (entity id, field, old value, new value) before writing, enabling
        human review.
  \item \textbf{Reject destructive deletions.} If a source item is deleted or removed
        from scope, the re-sync must flag it as a warning rather than silently removing
        the IR entity. A human or agent must confirm removal.
\end{enumerate}

\subsubsection{Stable IDs and Git-Friendly YAML}
\label{sec:stable-ids}

The combination of stable kebab-case IDs and line-oriented YAML serialisation produces
IR documents that diff well under version control:

\begin{lstlisting}[caption={Git diff of an IR update (status change + progress update)}]
  activities:
    - id: api-gateway-migration
      label: API Gateway Migration
      track: platform
      span: 2026-Q2
-     status: in-progress
-     progress: 0.4
+     status: done
+     progress: 1.0
      metadata:
        source: ado
        ado_id: 1024
\end{lstlisting}

Key properties of the YAML serialisation that preserve git-friendliness:

\begin{itemize}
  \item Fields within each entity are emitted in a \textbf{canonical order} (required
        fields first, optional fields in schema definition order). Reordering fields
        does not change semantics but does generate spurious diffs; canonical order
        prevents this.
  \item \textbf{No auto-generated IDs or timestamps} in non-\texttt{metadata} fields.
        The IR ID is derived from content and frozen; \texttt{ingested\_at} in
        \texttt{metadata} only changes when the entity is re-ingested from source.
  \item \textbf{Consistent quoting conventions.} Strings that are safe as bare YAML
        scalars are not quoted; strings with special characters use double quotes.
        Serialisers must apply this consistently to avoid quote-flip diffs.
  \item \textbf{Lists maintain insertion order}, not alphabetical. New entities appended
        to the end of a list produce a single-line diff at the bottom.
\end{itemize}

\subsubsection{Human Edits and Incremental Refinement}

A human author may open the IR YAML in any editor, add annotations, adjust labels,
add a milestone the ingester missed, or change a \texttt{status} value to reflect a
decision made after ingestion. These edits are valid IR operations and must survive
re-sync (per the strategy above).

An agent may similarly refine an IR: for example, an editorial agent might re-read the
document, notice that five activities in one track are unlabelled, and add descriptions.
Because the agent is editing a YAML file with stable IDs and a well-defined schema, this
is a safe, auditable operation. The MCP \texttt{validate\_timeline} tool can be called
after any edit to confirm the document remains valid.

The editing workflow is:

\begin{lstlisting}[caption={Iterative refinement loop (human or agent)}]
1. ingest(source, prompt)   -> roadmap.yaml          # initial ingestion
2. validate(roadmap.yaml)   -> valid, 0 errors        # confirm valid
3. edit(roadmap.yaml)       -> roadmap.yaml (modified) # human/agent edit
4. validate(roadmap.yaml)   -> valid, 0 errors        # re-confirm
5. render(roadmap.yaml, svg) -> roadmap.svg            # deterministic output
6. (commit roadmap.yaml + roadmap.svg to git)
\end{lstlisting}

Step~6 commits both the IR source and the rendered output. Because rendering is
deterministic, the SVG in the repository is always reproducible from the YAML. The YAML
is the source of truth; the SVG is a derived artefact that can be regenerated on demand.


%%============================================================================
\subsection{IR Gaps and Open Questions}
\label{sec:ir-gaps}
%%============================================================================

During the design of this section, three gaps or ambiguities in the IR contract were
identified. These are flagged here for Leslie (reviewer) and Mark to reconcile; they
have \emph{not} been silently resolved in this section's examples.

\begin{enumerate}
  \item \textbf{\texttt{today} field missing from metadata.} Leslie's scope decisions
        specify that the \texttt{now} symbolic date must be evaluated from an explicit
        \texttt{today} field in the IR (not from the system clock), in order to preserve
        determinism. However, the \texttt{metadata} table in Mark's IR contract does not
        include a \texttt{today} field. Until this is resolved, the annotation type
        \texttt{today-marker} with \texttt{date: now} is non-deterministic. Suggested
        resolution: add \texttt{today: date?} to the \texttt{metadata} schema; if absent,
        render-time clock is used and a warning is emitted.

  \item \textbf{Track sort field: \texttt{index} vs.\ \texttt{order}.} Mark's binding
        contract (well-formedness invariant~14) refers to a \texttt{track.index} field
        with the rule ``track index order values unique and non-negative''. The
        \texttt{04-ir.tex} specification defines the same field as \texttt{order}
        (type \texttt{int?}). The canonical field name should be resolved to one term;
        this section uses \texttt{order} to match the spec, but the contract should
        be updated to match.

  \item \textbf{\texttt{span} vs.\ \texttt{start}+\texttt{end} mutual exclusion.}
        The IR allows either \texttt{span} (shorthand) or \texttt{start} plus optional
        \texttt{end}, but does not specify the behaviour when both are present (e.g.,
        \texttt{span: 2026-Q2} and \texttt{start: 2026-05-01} coexist in the same
        activity). Should this be a schema error, or should one take precedence?
        For agent generation, structured output can prevent this from occurring, but the
        validation specification needs an explicit ruling. Suggested resolution: treat
        co-presence of \texttt{span} and \texttt{start}/\texttt{end} as a schema error.
\end{enumerate}
